\chapter*{Historical Note}
\phantomsection
\addcontentsline{toc}{chapter}{Historical Note}

(N.-B. --- The Roman numerals refer to the bibliography placed at the end of this note.)

The theory of quadratic forms, in its modern aspect, goes back hardly beyond the second half of the eighteenth \(^{e}\) century, and, as we shall see, it developed above all to meet the needs of Arithmetic, Analysis, and Mechanics. But the fundamental notions of this theory in reality made their appearance from the beginnings of "Euclidean" geometry, of which they form the framework. For this reason, one cannot trace its history without speaking, at least in a summary way, of the development of "elementary geometry" since antiquity. Of course, we can only concern ourselves with the evolution of a few general ideas, and the reader should not expect to find here precise information on the history of any particular theorem, concerning which it will suffice for us to refer to the specialized historical or didactic works (*). It goes without saying also that, when we speak below of the various possible interpretations of one and the same theorem in various algebraic or geometric languages, we in no way mean to say that these "translations" have at all times been as familiar as they are today; quite the contrary, it is the principal goal of this Note to show how, very gradually, mathematicians became aware of these affinities between questions often of very different aspect; we shall also have to show how, in doing so, they were led to introduce some coherence into the mass of geometric theorems bequeathed by the ancients, and finally to try to delimit exactly what should be understood by "geometry."

Apart from the Babylonians' discovery of the formula for solving the quadratic equation ((I), pp. 183-189), it is thus under their geometric disguise that the birth of the principal concepts of the theory of quadratic forms must be noted. These present themselves first as squares of distances (in the plane or in three-dimensional space), and the corresponding notion of "orthogonality" is introduced by means of the right angle, defined by Euclid as half of a straight angle (Elements, Book I, Def. 10); the notions of distance and right angle being related by the Pythagorean theorem, the keystone of the Euclidean edifice (*). The idea of angle seems to have been introduced very early in Greek mathematics (which probably received it from the Babylonians, well versed in the use of angles through their long astronomical experience). It is known that in the classical period, only angles less than two right angles are defined (Euclid's "definition" is moreover as vague and unusable as the one he gives for the straight line or the plane); the notion of orientation is not brought out, although Euclid uses (without axiom or definition) the fact that a line divides the plane into two regions, which he carefully distinguishes when necessary (**). At this stage, the idea of the group of plane rotations thus emerges only in a very imperfect manner, through the addition (introduced, also without explanation, by Euclid) of non-oriented angles of half-lines, which is defined, in principle, only when the sum is at most equal to two right angles (***). As for trigonometry, it is disdained by geometers and left to surveyors and astronomers; it is the latter (Aristarchus, Hipparchus, and especially Ptolemy (V)) who establish the fundamental relations between sides and angles of a right triangle (plane or spherical) and draw up the first tables (these are tables giving the chord of the arc cut off by an angle \(\theta < \pi\) on a circle of radius \(r\) , in other words the number \(2r \sin \frac{\theta}{2}\) ; the introduction of the sine, which is more convenient to handle, is due to the Hindu mathematicians of the Middle Ages); in the computation of these tables, the addition formula for arcs, unknown at that time, is replaced by the equivalent use of Ptolemy's theorem (possibly going back to Hipparchus) on quadrilaterals inscribed in a circle (cf. Esp. vect. top., chap. V, § 1, exerc. 5). It should also be noted that Euclid and Heron give propositions equivalent to the formula

\[
a ^ {2} = b ^ {2} + c ^ {2} - 2 b c \cos A
\]

between sides and angles of an arbitrary plane triangle; but one can hardly see in this a first appearance of the notion of the bilinear form associated with a metric form, for lack of the idea of a vector calculus that would emerge only in the nineteenth \(^{e}\) century.

Displacements (or motions, the distinction between the two notions not being clear in antiquity---nor even much later) were known to Euclid; but, for reasons unknown to us, he seems to show a marked reluctance to make use of them (for example in the "cases of equality of triangles," where one has the impression that he employs the notion of displacement only for want of having been able to formulate an appropriate axiom ((III bis), t. I, p. 225-227 and 249)); however, it is to the notion of displacement (rotation about an axis) that he has recourse for the definition of cones of revolution and spheres (Elements, Book XI, defs. 14 and 18), as does Archimedes for that of quadrics of revolution. But the general idea of transformation, applied to all of space, is almost foreign to mathematical thought before the end of the eighteenth \(^{e}\) century (*);

and before the seventeenth \(^{e}\) century, there is likewise no trace of the notion of composition of motions, nor a fortiori of composition of displacements. This does not mean, of course, that the Greeks were not particularly sensitive to the “regularities” and “symmetries” of figures, which we now connect with the notion of the group of displacements; their theory of regular polygons and even more that of regular polyhedra — one of the most remarkable chapters of all their mathematics — is there to prove the contrary (*).

Finally, the last of the essential contributions of Greek mathematics, in the domain that concerns us, is the theory of conics (as for quadrics, the Greeks knew only certain quadrics of revolution, and did not carry their study very far, except for the sphere). It is interesting to note here that, although the Greeks never had the idea of the fundamental principle of analytic geometry (essentially for lack of a manageable algebra), they commonly used, for the study of particular "figures," the "ordinates" with respect to two (or even more than two) axes in the plane (closely related to the figure, which is one of the fundamental points where their method differs from that of Fermat and Descartes, whose axes are fixed independently of the figure under consideration). In particular, the first examples of conics (other than the circle) that appear in connection with the problem of the duplication of the cube are the curves given by the equations \(y^{2} = ax\) , \(y = bx^{2}\) , \(xy = c\) (Menaechmus, a student of Eudoxus, mid- \(1v^{e}\) century) (**); and it is the equation of conics (usually with respect to two oblique axes formed by a diameter and the tangent at one of its points of intersection with the curve) that is most often used in the study of problems relating to these curves (whereas "focal" properties play only a very subdued role, contrary to what school traditions dating back only to the \(x1x^{e}\) century might lead one to believe). From this vast theory, what we must above all retain here is the notion of conjugate diameters (already known to Archimedes), and the property that now serves as the definition of the polar of a point, given by Apollonius (IV) when the point is exterior to the conic (the polar thus being for him the line joining the points of contact of the tangents drawn from this point); from our point of view, these are two examples of "orthogonality" with respect to a quadratic form distinct from the metric form, but of course the connection between these notions and the classical notion of perpendiculars could not possibly have been conceived at that time.

There is scarcely any other progress to report before Descartes and Fermat; but from the very beginnings of analytic geometry, the algebraic theory of quadratic forms begins to emerge from its geometric shell: Fermat knows that a second-degree equation in the plane represents a conic ((VII a), p. 100-102) and sketches analogous ideas on quadrics (VII b). With the development of analytic geometry in 2 and 3 dimensions during the eighteenth \(^{e}\) century there appear (especially in connection with conics and quadrics) two of the central problems of the theory: the reduction of a quadratic form to a sum of squares and the search for its "axes" with respect to the metric form. For conics, these two problems are too elementary to give rise to important algebraic progress; for an arbitrary number of variables, the first is solved by Lagrange in 1759, in connection with the maxima of functions of several variables (IX a). But this problem is almost immediately overshadowed by that of the search for axes, even before the invariance of the rank had been formulated (*); as for the law of inertia, it is discovered only around 1850 by Jacobi (XIX b), who proves it by the same reasoning as is used today, and Sylvester (XX), who merely states it as almost evident (**).

The problem of reducing a quadric to its axes already presents substantially greater algebraic difficulties than the analogous problem for conics; and Euler, who is the first to tackle it, is not in a position to prove the reality of the eigenvalues, which he admits after a sketch of a justification lacking probative value ((VIII a), p. 379-392) (**). If this point is correctly established around 1800, one must wait for Cauchy to prove the corresponding theorem for forms in an arbitrary number \(n\) of variables (XIV b). It is also Cauchy who, around the same time, proves that the characteristic equation giving the eigenvalues is invariant under any change of rectangular axes ((XIV a), p. 252) (****); but for \(n = 2\) or \(n = 3\) , this invariance was intuitively "evident" because of the geometric interpretation of the eigenvalues by means of the axes of the corresponding conic or quadric. Moreover, in the course of research on this subject, the elementary symmetric functions of the eigenvalues also presented themselves in a natural way (with various geometric interpretations, relating in particular to Apollonius's theorems on conjugate diameters), and in particular the discriminant, which (known for a long time for n = 2 in connection with the theory of the second-degree equation) appears for the first time for n = 3 in Euler ((VIII a) p. 382); the latter encounters it in connection with the classification of quadrics (in expressing the condition for a quadric to have no point at infinity) and does not mention its invariance under changes of rectangular axes. But a little later, with the beginnings of the arithmetic theory of quadratic forms with integer coefficients, Lagrange notes (for n = 2) a particular case of invariance of the discriminant under a linear but non-orthogonal change of variables ((IX b), p. 699), and Gauss establishes, for n = 3, the "covariance" of the discriminant under any linear transformation ((XI a), p. 301-302) (*). Once the general formula for the multiplication of determinants had been proved by Cauchy and Binet, the extension of Gauss's formula to an arbitrary number of variables was immediate; it is this that, around 1845, was to give the first impetus to the general theory of invariants.

To the two notions which, among the Greeks, took the place of the theory of displacements --- namely the notion of motion and that of "symmetry" of a figure --- there is added a third in the seventeenth \(^{e}\) and eighteenth \(^{e}\) centuries with the problem of the change of rectangular axes, which is substantially equivalent to this theory. Euler devotes several works to this question, striving above all to obtain manageable parametric representations for the formulas of change of axes. We know what use Mechanics was to make of the three angles that he introduces for this purpose for \(n = 3\) ((VIII a), p. 371-378). But he does not stop there; he considers in 1770 the general problem of orthogonal transformations for arbitrary \(n\) , observes that one arrives here at the goal by introducing \(n(n - 1)/2\) angles as parameters, and finally, for \(n = 3\) and \(n = 4\) , gives rational representations for rotations (in terms, respectively, of 4 homogeneous parameters and of 8 homogeneous parameters linked by a relation), which are none other than those obtained later by means of the theory of quaternions (cf. § 9, exerc. 15 and 16), and whose origin he does not indicate (VIII c) (**).

On the other hand, Euler also indicates how to translate analytically the search for the “symmetries” of plane figures, and it is in this connection that he is led to prove, in substance, that a plane displacement is a rotation, or a translation, or a translation followed by a symmetry ((VIII a), pp. 197-199). The rise of Mechanics at that time moreover leads to the general study of displacements; but at first only “infinitely small” displacements tangent to continuous motions are in question: these are apparently the only ones that occur in the researches of Torricelli, Roberval, and Descartes on the composition of motions and the instantaneous center of rotation for plane motions (cf. Hist. Note of Book IV, chaps. I-II-III). The latter is defined in a general way by Johann Bernoulli; d'Alembert in 1749, Euler the following year, extend this notion by proving the existence of an instantaneous axis of rotation for motions leaving a fixed point. The analogous theorem for finite displacements is stated only in 1775 by Euler (VIII d), in a memoir in which he discovers at the same time that the determinant of a rotation is equal to 1; the following year, he proves the existence of a fixed point for plane similarities (VIII e). But it will be necessary to wait for the works of Chasles, starting in 1830 (XV a), to finally have a coherent theory of finite and infinitely small displacements.
% semantic-fragments: legacy-input-seam
\relax{}
\begin{center}
\(\ast\qquad\ast\)
\end{center}

We thus arrive at what may be called the golden age of geometry, which falls roughly between the publication dates of Monge's Descriptive Geometry (1795) (X) and F. Klein's "Erlangen Program" (1872) (XXV b). The essential advances we owe to this sudden revival of geometry are the following:

\begin{enumerate}
\def\labelenumi{\Alph{enumi})}
\item
  The notion of an element at infinity (point, line, or plane), introduced by Desargues in the seventeenth century (VI), but which hardly manifests itself in the eighteenth century only as an abuse of language, is rehabilitated and systematically used by Poncelet (XII), who thus makes projective space the general framework for all geometric phenomena.
\item
  At the same time, with Monge, and above all Poncelet, the transition to complex projective geometry takes place. The notion of imaginary point, sporadically used during the eighteenth century, is exploited here (concurrently with that of the point at infinity) to yield statements independent of the "cases of figure" of real affine geometry. If, at first, the justifications offered in support of these innovations remain highly awkward (especially on the part of the adherents of the "synthetic" school of geometry, where the use of coordinates came to be regarded as a stain), one cannot fail to recognize therein, under the name of "principle of contingent relations" in Monge, or of "principle of continuity" in Poncelet, the first germ of the idea of "specialization" in modern algebraic geometry (*).
\end{enumerate}

One of the first results arising from these conceptions is the observation that, in complex projective space, all nondegenerate conics (resp. quadrics) are of the same nature; this leads Poncelet to the discovery of the "isotropic" elements: "Circles placed arbitrarily in a plane," he says, "are therefore not entirely independent of one another, as one might at first believe; they have ideally two imaginary points in common at infinity" ((XII), p.~48). Further on, he likewise introduces the "umbilic," the imaginary conic at infinity common to all spheres ((XII), p.~370); and if he does not speak particularly of the isotropic generators of the sphere, he at least explicitly emphasizes the existence of rectilinear generators, real or imaginary, for all quadrics (ibid., p.~371) (*); notions of which his successors (notably Plücker and Chasles), even more than he himself, make great use, in particular in the study of the "focal" properties of conics and quadrics.

\begin{enumerate}
\def\labelenumi{\Alph{enumi})}
\setcounter{enumi}{2}
\item
  The notions of point transformation and composition of transformations are likewise formulated in a general way and introduced systematically as means of proof. Apart from displacements and projections, only a few particular transformations had been known up to that time: certain planar projective transformations, of the type \(x' = a/x\) , \(y' = y/x\) used by La Hire and Newton, the “affinity” \(x' = ax\) , \(y' = by\) from Clairaut and Euler, and finally certain particular quadratic transformations, in Newton again, Maclaurin, and Braikenridge. Monge, in his Descriptive Geometry, shows all the use that can be made of plane projections in three-dimensional geometry. In Poncelet, one of the systematic methods of proof, employed to excess, consists in reducing by projection the properties of conics to those of the circle (a method already applied on occasion by Desargues and Pascal); and in order to pass similarly from a quadric to a sphere, he invents the first example of a projective transformation in space, the "homology" ((XII), p.~357); finally, it is also he who introduces the first examples of birational transformations of a curve into itself. In 1827, Möbius ((XIII a), p.~217) (and independently Chasles in 1830 (XV b)) define the most general projective linear transformations; at the same period there appear inversion (cf.~§ 10, exerc. 13) and other types of quadratic transformations, whose study will inaugurate the theory of birational transformations, which will develop in the second half of the nineteenth century.
\item
  The notion of duality appears in full light and is consciously linked to the theory of bilinear forms. The theory of poles and polars with respect to conics, which since Apollonius had made progress only in the work of Desargues and La Hire, is extended to quadrics by Monge, who, together with his students, perceives the possibility of transforming known theorems into new results by this means (*). But it is again to Poncelet that the credit belongs for having erected these remarks into a general method in his theory of transformations "by reciprocal polars," and for having made of it a particularly effective tool of discovery. A little later, notably with Gergonne, Plücker, Möbius, and Chasles, the general notion of duality detaches itself from the connection with quadratic forms, still too narrow in Poncelet. In particular, Möbius, in examining the various possibilities of duality in 3-dimensional space (defined by a bilinear form), discovers in 1833 the duality with respect to an alternating bilinear form (XIII b) (**), especially studied in the nineteenth century, in the form of the theory of “linear complexes” (cf.~§ 10, exerc. 16) and developed in connection with the “geometry of lines” and the “Plücker coordinates” introduced by Cayley, Grassmann, and Plücker around 1860.
\item
  From the very beginnings of projective geometry, the intensive study of the properties of classical geometry in their relations with projective space had quickly led to dividing them into "projective properties" and "metric properties"; and it is doubtless no exaggeration to see in this separation one of the clearest manifestations, at that time, of what was to become the modern notion of structure. But Poncelet, who first introduced this distinction and this terminology, was already aware of what connects these two types of properties; and, addressing in his Treatise the problems concerning angles, whose properties "do not seem to form part of those we have called projective\ldots, they nevertheless follow in so simple a manner," he said, "from the principles that form the basis {[}of this work{]}\ldots, which I do not believe that any other geometric theory can lead to in a manner at once more direct and simpler. One will not be at all surprised by this, if one considers that the projective properties of figures are necessarily the most general of those that can belong to them; so that they must include, as simple corollaries, all the other particular properties or relations of extension" ((XII), p.~248). To tell the truth, after this declaration, one is somewhat surprised to see him approach questions of angles in a very roundabout way, by attaching them to the focal properties of conics, instead of bringing the cyclic points directly into play; and in fact, it was only 30 years later that Laguerre (still a student at the École Polytechnique) gave the expression of an angle between lines by means of the cross-ratio of these lines and the isotropic lines with the same origin (cf.~§ 10, exerc. 5) ((XXI), vol. II, p.~13). Finally, with Cayley (XVIII d) there is clearly expressed the fundamental idea that the "metric" properties of a plane figure are none other than the "projective" properties of the figure augmented by the cyclic points --- a decisive milestone toward the "Erlangen program".
\item
  Hyperbolic non-Euclidean geometry, which came into being around 1830, at first remains somewhat apart from the movement whose main lines we are tracing. Arising from concerns of an essentially logical order touching on the foundations of classical geometry, this new geometry is presented by its inventors (*) in the same axiomatic and "synthetic" form as Euclid's geometry, and without any connection to projective geometry (whose introduction following the classical model even seemed excluded a priori, since the notion of a unique parallel disappears in this geometry); it is doubtless for this reason that for a long time it attracted little interest from the French, German, and English schools of projective geometry. Thus, when Cayley, in the fundamental memoir cited above (XVIII d), had the idea of replacing the cyclic points (considered as a conic "degenerate tangentially") by an arbitrary conic (which he called the "absolute"), he in no way thought of connecting this idea to the geometry of Lobachevsky-Bolyai, although he indicates how his conception leads to new expressions for the "distance" between two points, and although he mentions its links with spherical geometry. The situation changes around 1870, when the non-Euclidean geometries, following the diffusion of the works of Lobachevsky, and the publication of the works of Gauss and of Riemann's inaugural lecture, came to the forefront of mathematical actuality. Following the path traced by Riemann, Beltrami, without knowing Cayley's work, rediscovered in 1868 the expressions for distance given by the latter, but in a quite different context, by considering the interior of a circle as an image of a surface of constant curvature, in which geodesics are represented by straight lines (XXIV); it is Klein who, two years later, made (independently of Beltrami) the synthesis of these various points of view, which he completed by the discovery of elliptic non-Euclidean space (XXV a) (**).
\item
  In the second half of the period we are considering here, a period of critical reflection sets in, during which the partisans of "synthetic" geometry, not content with having banished coordinates from their proofs, claim to do without real numbers even in the axioms of geometry. The principal representative of this school is von Staudt, who essentially succeeded in performing this tour de force (XXIII), much admired in his time and even well into the twentieth century; and if today one no longer attributes the same importance to ideas of this order, whose possibilities for fruitful application have proved rather meager, one must nevertheless recognize that the efforts of von Staudt and his disciples contributed to clarifying ideas on the role of real or complex "scalars" in classical geometry, and thereby to introducing the modern conception of geometries over an arbitrary base field.
\end{enumerate}

Around 1860, "synthetic" geometry was at its zenith, but the end of its reign was fast approaching. Having remained heavy and inelegant throughout the eighteenth century, analytic geometry, in the hands of Lamé, Bobillier, Cauchy, Plücker, and Möbius, finally acquired the elegance and concision that would enable it to compete on equal terms with its rival. Above all, from about 1850 onward, the ideas of group and invariant, at last formulated precisely, gradually took center stage, and it became apparent that the theorems of classical geometry are nothing other than the expression of identical relations among invariants or covariants of the similarity group (*), just as those of projective geometry express the identities (or "syzygies") among covariants of the projective group. This is the thesis masterfully expounded by F. Klein in the celebrated "Erlangen program" (XXV b), where he advocates abandoning the sterile controversies between the "synthetic" and "analytic" tendencies; if, he says, the accusation leveled against the latter of granting a privileged role to an arbitrary system of axes "was only too often justified with regard to the defective manner in which the coordinate method was formerly used, it collapses when it comes to a rational application of this method.\ldots{} The domain of spatial intuition is not forbidden to the analytic method.\ldots{}", and he emphasizes that “one must not underestimate the advantage that a well-suited formalism brings to subsequent research, in that it anticipates thought, so to speak” ((XXV b), pp. 488–490).

This leads to a rational and "structural" classification of the theorems of "geometry" according to the group from which they derive: the linear group for projective geometry, the orthogonal group for metric questions, and the symplectic group for the geometry of the "linear complex." But under this pitiless clarity, classical geometry—with the exceptions of algebraic geometry and differential geometry (**), now constituted as autonomous sciences—suddenly withers and loses all its brilliance. Already the generalization of methods based on the use of transformations had rendered the formation of new theorems somewhat mechanical: "Today," said Chasles in 1837 in his \emph{Aperçu historique}, "anyone may come forward, take any known truth, and subject it to the various general principles of transformation; he will extract from it other truths, different or more general; and these will be susceptible to similar operations; so that one may multiply, almost to infinity, the number of new truths deduced from the first.\ldots{} Anyone who wishes may therefore, in the current state of science, generalize and create in Geometry; genius is no longer indispensable for adding a stone to the edifice" ((XV b), p.~268-269). But the situation becomes much clearer with the progress of invariant theory, which finally succeeds (at least for the “classical” groups) in formulating general methods that make it possible in principle to write out all algebraic covariants and all their “syzygies” in a purely automatic manner; a victory which, at the same time, marks the death, as a field of research, of classical invariant theory itself, and of “elementary” geometry, which has practically become a mere dictionary for it. To be sure, nothing allows one to foresee a priori, among the infinity of “theorems” that can thus be unrolled at will, which ones will have, in an appropriate geometric language, a simplicity and elegance comparable to the classical results, and there remains a restricted domain where numerous amateurs continue to exercise themselves happily (geometry of the triangle, of the tetrahedron, of algebraic curves and surfaces of low degree, etc.). But for the professional mathematician, the mine is exhausted, since there are no longer any problems of structure there, capable of having repercussions on other parts of mathematics; and this chapter of group theory and invariant theory may be considered closed until further notice. (*)

\[
\ast \ast
\]

Thus, after the Erlangen program, Euclidean and non-Euclidean geometries, from a purely algebraic point of view, became simple languages, more or less convenient, for expressing the results of the theory of bilinear forms, whose progress goes hand in hand with that of the theory of invariants (*). Everything concerning the notion of the rank of a bilinear form and the relations between these forms and linear transformations is definitively clarified by the work of Frobenius (XXVII a). It is also to Frobenius that we owe the canonical expression of an alternating form on a free Z-module (§ 5, no. 1, th. 1) (XXVII b); however, skew-symmetric determinants had already appeared in Pfaff, at the beginning of the century, in connection with the reduction of differential forms to a normal form; Jacobi, who in 1827 takes up this problem again (XIX a), knows that a skew-symmetric determinant of odd order is zero, and it is he who forms the expression of the Pfaffian and shows that it is a factor of the skew-symmetric determinant of even order; but he had not perceived that the latter is the square of the Pfaffian, and this point was established only by Cayley in 1849 (XVIII b). The notion of the symmetric bilinear form associated with a quadratic form is the most elementary case of the process of "polarization," one of the fundamental tools of the theory of invariants. Under the name of "scalar product," this notion would enjoy immense fortune, first with the popularizers of "vector calculus," then, from the twentieth century, thanks to the unforeseen generalization brought to it by the theory of Hilbert space (see the historical note in Book V). It is also this latter theory that will bring to light the notion of the adjoint of an operator (which previously had scarcely manifested itself except in the theory of linear differential equations, and, in tensor calculus, through the dance of co- and contravariant indices under the baton of the metric tensor); it is this theory, finally, that will give full prominence to the notion of Hermitian form, introduced first by Hermite in 1853 in connection with arithmetic investigations ((XXII), p.~237), but which remained somewhat on the margins of the main mathematical currents until around 1925 and the applications of complex Hilbert spaces to quantum theories.

The study of the orthogonal group and the group of similarities --- clearly conceived and treated as such since the middle of the nineteenth century, and having become the heart of the theory of quadratic forms — as well as of the other "classical" groups (the linear group, the symplectic group, and the unitary group) — takes on, on the other hand, an increasingly important role. We can only mention here the essential part played by these groups in the theory of Lie groups and differential geometry on the one hand, and in the arithmetic theory of quadratic forms (see, for example, (XXXIII) and (XXXV)) on the other (**); to this circumstance, as well as to the extension of the concept of duality to the most diverse questions, is due the fact that there is scarcely any modern mathematical theory in which bilinear forms do not intervene in one way or another. We must in any case note that it was the study of the rotation group (in three dimensions) that led Hamilton to the discovery of quaternions (XVII); this discovery was generalized by W. Clifford, who in 1876 introduced the algebras bearing his name and proved that they are tensor products of quaternion algebras, or of quaternion algebras and a quadratic extension (XXVIII). Rediscovered four years later by Lipschitz (XXIX), who used them to give a parametric representation of orthogonal transformations in n variables (generalizing those that Cayley had obtained for n = 3 (XVIII a) and n = 4 (XVIII c) via the theory of quaternions (cf. § 9, exercises 15 and 16)), these algebras, and the notion of "spinor" derived from them (see (XXXII) and (XXXIV)), were also to enjoy great popularity in the modern era by virtue of their use in quantum theories.

Finally, we have only a word left to say about the evolution of ideas that led to the almost total abandonment of any restriction on the ring of scalars in the theory of sesquilinear forms — a tendency common to all modern algebra, but one that perhaps manifested itself here earlier than elsewhere. We have already noted the fruitful introduction of geometry over the field of complex numbers (which, moreover, throughout the nineteenth century, was not without a perpetual and sometimes perilous confusion between this geometry and real geometry); here the clarity comes above all from the axiomatic studies of the end of the nineteenth century on the foundations of geometry (XXX). In the course of these investigations, Hilbert and his followers, in particular, while examining the relations among the various axioms, were led to construct suitable counterexamples, where the “base field” (commutative or not) possessed more or less pathological properties, and they thus accustomed mathematicians to “geometries” of an entirely new type. From the analytic point of view, Galois had already considered linear transformations where coefficients and variables took their values in a finite prime field ((XVI), p.~27); in developing these ideas, Jordan (XXVI) is naturally led to envisage the classical groups over these fields, groups whose intervention manifests itself in varied domains of mathematics. Dickson, around 1900, extended Jordan's investigations to all finite fields, and more recently it has been realized that a large part of the Jordan–Dickson theory extends to the case of an absolutely arbitrary “base field”; this is due essentially to the general properties of isotropic vectors and to Witt's theorem, which, trivial in the classical cases, were established for an arbitrary base field only in 1936 (XXXI) (*).

But in thus pushing the study of sesquilinear forms toward an ever greater “abstraction,” it proved extremely suggestive to retain as is the terminology which, in the case of 2- and 3-dimensional spaces, came from classical geometry, and to extend it to the n-dimensional case and even to infinite-dimensional spaces. Outmoded as an autonomous and living science, classical geometry was thereby transfigured into a universal language of contemporary mathematics, of incomparable flexibility and convenience.

